
\subsubsection{2.1 Reward Misspecification and RLHF Overoptimization}

Reinforcement Learning from Human Feedback [Christiano et al. 2017; Ouyang et al. 2022] trains reward models on human preference labels and optimizes language model policies against these reward models. A central concern is reward misspecification: reward models learn proxy signals that diverge from genuine quality, leading to Goodhart's Law violations as optimization pressure increases [Gao et al. 2022].

Pan et al. [2022] provide a systematic analysis of reward misspecification, demonstrating that even small errors in the reward model can compound over RLHF training to produce measurable accuracy drops on held-out benchmarks. Coste et al. [2023] show that annotation variance — disagreement between annotators on the same prompt — limits reward model reliability and contributes to overoptimization. Both lines of work treat the annotation signal as stationary noise rather than as a potentially time-varying source of bias. The present work examines whether the annotation signal itself shifts direction across annotation rounds.

\subsubsection{2.2 RLHF Annotation Dynamics and Dataset Structure}

Bai et al. [2022] describe the Anthropic HH-RLHF dataset as a product of sequential annotation phases, with annotators evaluating pairs of AI-generated responses. Stiennon et al. [2020] describe the OpenAI WebGPT comparisons dataset, collected via a crowdwork design intended to enable within-annotator analysis. Neither paper analyzes whether annotator preferences shift directionally across annotation phases or sessions.

Ziegler et al. [2019] study the effect of reward model scale on RLHF quality under the assumption of preference stationarity. The temporal structure of multi-round annotation datasets — which in principle enables preference stationarity analysis — has not previously been exploited for this purpose.

\subsubsection{2.3 LLM-as-Judge and Evaluation Adaptation}

Thakur and Kambhampati [2024] document criterion shift in LLM-as-judge evaluation: language model evaluators adapt their assessment criteria under AI text exposure. The present work examines an analogous phenomenon in human annotators contributing to RLHF training data, and connects it to a specific causal pathway involving stylistic coefficient drift.

Liang et al. [2023] and related work on verbosity bias in LLM evaluation document that AI evaluators tend to prefer longer, more structured responses independent of quality. The present findings on verbosity coefficient drift in human annotation are consistent with this pattern.

\subsubsection{2.4 Automation Bias}

Automation bias — the tendency of humans to defer to automated outputs under decision uncertainty — is documented across aviation, medicine, and decision support [Skitka et al. 1999; Parasuraman \& Manzey 2010]. Lee \& See [2004] and Dietvorst et al. [2015] document criterion drift and reversal of algorithmic aversion under repeated exposure. These findings motivate the hypothesis that RLHF annotators, under repeated exposure to AI-generated text, may internalize AI stylistic features as proxies for quality.

\subsubsection{2.5 Positioning}

This work occupies a gap defined by the intersection of {training data annotation} and {measuring directional drift}. Prior work addresses downstream symptoms of annotation drift — reward misspecification, benchmark degradation, label noise — without tracing these to a temporal source in the annotation process. The AAI is introduced as a scalar, pre-registerable instrument for this upstream measurement.

